% Section 3.1. Worked example: data. Prompts and row provenance are in Supplement S1 and S3.

\section{Worked Example}
\label{sec:data}

\subsection{Data}

The worked example uses the corpus of \citet{keough2026plausible}, in which language models answered the eight-item Patient Health Questionnaire depression scale (PHQ-8; \citealp{kroenke2009phq8}) as demographic personas.
Four models, GPT-4o-mini, Gemini-3-Flash, DeepSeek-V3 and GLM-4.7, generated the corpus through OpenRouter between 28 December 2025 and 11 January 2026, at the provider's default decoding. % R: corpus_v3_provenance.csv
The persona grid crosses race (Asian, Black, Hispanic, multiracial, White), gender identity (cisgender and transgender women and men), income (three levels, stated as a dollar-and-insurance label from ``Low (<\$35k, Medicaid)'' to ``High (>\$250k, Concierge)'') and relationship status (married, single), for 120 personas.
Each persona answered 30 times under each of two framings, one listing its attributes as fields (clinical) and one giving them as a first-person self-description (narrative).
A shared system prompt cast the model as a ``Clinical Simulation Engine'' answering from ``the statistical likelihood of symptoms for this specific demographic intersection in the US population''.
The corpus holds 28,800 assessments, of which 28,799 have a valid PHQ-8, and Supplements S1 and S3 give the prompts and the provenance of every row. % R: corpus_v3_provenance.csv; l1_personfit_main.csv

The population reference is the 36,274 adults in NHANES 2005--2018 who answered all eight PHQ-8 items, analyzed with the survey's weights and design \citep{johnson2013nhanes,chen2018nhanes,nchs2018analytic}. % R: l1_nhanes_sample_flow.csv; analysis_brm/L4.md
Income is the ratio of family income to the poverty line, banded below 1.3, from 1.3 to 3.5, and above 3.5, and each persona's income label is mapped to one band.
Forty-eight of the 120 personas have a matching group in the survey.
NHANES records sex rather than gender identity and does not identify multiracial adults on their own, and transgender and multiracial personas therefore enter only levels 1 and 4, which read the whole simulated population.
NHANES 2021--2023 and other reference choices are sensitivity analyses in Supplement S3.

The NHANES standard deviation of the PHQ-8 total sets the gate tolerances at 0.98 and 0.49 points and the default level-3 tolerance at about 2 points. % R: gate_dstudy.csv
No NHANES subgroup by sex, race and ethnicity, income or survey cycle departs from 1 by more than a factor of 1.36 at level 1, or 1.27 at the near end of its interval, and level 1 uses $\tau = 1.5$. % R: l1_personfit_tolerance.csv
Level 2 reads seven contrasts for each model, the sex gap, three income gaps and three racial gaps, with Bonferroni-adjusted 98.6\% intervals.
Levels 2 and 3 use the mean of all 30 draws of each persona, more than the gate requires for any model, and level 4 tests invariance across sex, race and income.
